\documentclass{article}
\usepackage{amsmath,amssymb,amsthm}
\usepackage{algorithm}
\usepackage{algorithmic}
\usepackage{enumitem}
\usepackage{hyperref}
\newtheorem{theorem}{Theorem}
\newtheorem{lemma}{Lemma}
\newtheorem{assumption}{Assumption}
\newcommand{\grad}{\nabla}
\newcommand{\inner}[2]{\langle #1,#2\rangle}
\newcommand{\diam}{\mathrm{diam}}
\newcommand{\gap}{\mathcal{G}}
\begin{document}

%%%%%%%%%%%%%%%%%%%%%%%%%%%%%%%%%%%%%%%%%%%%%%%%%%%%%%%%%%%%
\section*{Algorithm Name}
\textbf{Frank–Wolfe (Conditional Gradient) for Non‑Convex Smooth Optimization}
%%%%%%%%%%%%%%%%%%%%%%%%%%%%%%%%%%%%%%%%%%%%%%%%%%%%%%%%%%%%
\section*{Mathematical Formulation}
Let $f:\mathbb{R}^{d}\rightarrow\mathbb{R}$ be differentiable and $C\subset\mathbb{R}^{d}$ a non‑empty compact convex set.  
Given a stepsize schedule $\{\gamma_k\}_{k\ge 0}\subset (0,1]$, the algorithm generates the sequence $\{x_k\}$ as follows:
\begin{align}
s_k &:= \arg\min_{s\in C}\;\inner{\grad f(x_k)}{s}, \tag{1}\\[4pt]
x_{k+1} &:= (1-\gamma_k)\,x_k + \gamma_k\, s_k . \tag{2}
\end{align}
The quantity
\[
\gap(x_k):=\max_{s\in C}\inner{\grad f(x_k)}{x_k-s}
      =\inner{\grad f(x_k)}{x_k-s_k}
\]
is the \emph{Frank–Wolfe gap} and serves as a first‑order stationarity measure.

%%%%%%%%%%%%%%%%%%%%%%%%%%%%%%%%%%%%%%%%%%%%%%%%%%%%%%%%%%%%
\section*{Pseudocode}
\begin{algorithm}[H]
\caption{Frank–Wolfe for Non‑Convex Smooth Problems}
\begin{algorithmic}[1]
\REQUIRE Initial point $x_0\in C$, stepsize rule $\{\gamma_k\}$, maximum iterations $K$.
\FOR{$k=0$ \TO $K-1$}
    \STATE Compute gradient $g_k \gets \grad f(x_k)$.
    \STATE \textbf{Linear Minimization Oracle (LMO)}:
    \[
        s_k \gets \arg\min_{s\in C}\inner{g_k}{s}.
    \]
    \STATE Choose stepsize $\gamma_k\in(0,1]$ 
    \COMMENT{e.g. $\gamma_k=\frac{2}{k+2}$ or exact line‑search.}
    \STATE Update $x_{k+1}\gets (1-\gamma_k)x_k+\gamma_k s_k$.
\ENDFOR
\RETURN $x_K$ (or the iterate with smallest gap).
\end{algorithmic}
\end{algorithm}
%%%%%%%%%%%%%%%%%%%%%%%%%%%%%%%%%%%%%%%%%%%%%%%%%%%%%%%%%%%%
\section*{Dry Run Example}
Consider the univariate quadratic 
\[
f(w)=(w-3)^2,\qquad C=[-5,5],\qquad x_0=0,\qquad \gamma_k\equiv 0.1 .
\]

\begin{center}
\begin{tabular}{c|c|c|c|c}
$k$ & $x_k$ & $\grad f(x_k)=2(x_k-3)$ & $s_k=\arg\min_{s\in C}\inner{\grad f(x_k)}{s}$ & $x_{k+1}=(1-\gamma_k)x_k+\gamma_k s_k$ \\ \hline
0 & $0$ & $-6$ & $s_0=5$ (largest $s$ because gradient is negative) & $x_1=0.9\cdot0+0.1\cdot5=0.5$ \\ \hline
1 & $0.5$ & $-5$ & $s_1=5$ & $x_2=0.9\cdot0.5+0.1\cdot5=0.95$ \\ \hline
2 & $0.95$ & $-4.1$ & $s_2=5$ & $x_3=0.9\cdot0.95+0.1\cdot5=1.355$ 
\end{tabular}
\end{center}
Objective values: $f(x_0)=9$, $f(x_1)=6.25$, $f(x_2)=4.2025$, $f(x_3)=2.887$.
The iterates move towards the minimiser $w^\star=3$ while respecting the feasible set $C$.

%%%%%%%%%%%%%%%%%%%%%%%%%%%%%%%%%%%%%%%%%%%%%%%%%%%%%%%%%%%%
\section*{Assumptions}
\begin{assumption}[Smoothness]\label{ass:smooth}
$f$ is $L$‑smooth on an open set containing $C$:
\[
\forall x,y\in C,\qquad
\|\grad f(x)-\grad f(y)\| \le L\|x-y\|.
\]
Equivalently (Descent Lemma),
\[
f(y)\le f(x)+\inner{\grad f(x)}{y-x}+ \frac{L}{2}\|y-x\|^2 . \tag{3}
\]
\end{assumption}
\begin{assumption}[Compactness]\label{ass:compact}
$C$ is convex, compact and has finite diameter
\[
D:=\diam(C)=\max_{u,v\in C}\|u-v\|<\infty .
\]
\end{assumption}
\begin{assumption}[Bounded Gradient on $C$]\label{ass:gradbound}
There exists $G>0$ such that $\|\grad f(x)\|\le G$ for all $x\in C$.
\end{assumption}
No convexity of $f$ is required; the algorithm is analysed for non‑convex $f$.

%%%%%%%%%%%%%%%%%%%%%%%%%%%%%%%%%%%%%%%%%%%%%%%%%%%%%%%%%%%%
\section*{Main Convergence Theorem}
\begin{theorem}[Convergence of Frank–Wolfe Gap]\label{thm:convergence}
Under Assumptions \ref{ass:smooth}–\ref{ass:gradbound}, let the stepsizes be
\[
\gamma_k = \min\!\Big\{1,\; \frac{2}{k+2}\Big\}.
\]
Then for any $K\ge 1$ the minimum Frank–Wolfe gap among the first $K$ iterates satisfies
\[
\min_{0\le k < K}\gap(x_k)\;\le\;
\frac{2\big(f(x_0)-f^\star\big)}{K} \;+\; \frac{L D^{2}}{K},\tag{4}
\]
where $f^\star:=\min_{x\in C}f(x)$. Consequently,
\[
\gap(x_{k}) = \mathcal O\!\Big(\frac{1}{\sqrt{k}}\Big)\quad\text{in the sense that}\quad
\frac{1}{K}\sum_{k=0}^{K-1}\gap(x_k)=\mathcal O\!\Big(\frac{1}{\sqrt{K}}\Big).
\]
\end{theorem}
The theorem shows that even without convexity the algorithm drives the Frank–Wolfe gap to zero at the same $\mathcal O(1/k)$ rate for the best iterate, and an $\mathcal O(1/\sqrt{k})$ rate in expectation over the iterates.

%%%%%%%%%%%%%%%%%%%%%%%%%%%%%%%%%%%%%%%%%%%%%%%%%%%%%%%%%%%%
\section*{Theoretical Properties}
\begin{itemize}[leftmargin=*]
\item \textbf{Stability}: The update is a convex combination of two points in $C$, hence $x_{k}\in C$ for all $k$ by induction.
\item \textbf{Hyper‑parameter Sensitivity}: The standard diminishing schedule $\gamma_k=2/(k+2)$ guarantees the bound (4). Larger constant stepsizes may still converge but the bound deteriorates; a line‑search that exactly minimizes $f\big((1-\gamma)x_k+\gamma s_k\big)$ yields at least as good a bound.
\item \textbf{Comparison to Gradient Descent}: Gradient descent requires projection onto $C$, which can be expensive. Frank–Wolfe replaces the projection by a linear minimization oracle, often cheaper for structured $C$ (e.g., simplex, trace‑norm ball). The convergence rate matches that of projected gradient descent for smooth non‑convex problems up to constant factors.
\item \textbf{Special Cases}:
    \begin{enumerate}
        \item If $f$ is convex, the bound (4) reduces to the classic $\mathcal O(1/k)$ rate for the primal gap $f(x_k)-f^\star$.
        \item If $C$ is a polytope and the LMO is exact, the iterates are sparse (Carathéodory’s theorem) – a structural advantage absent in projected methods.
    \end{enumerate}
\end{itemize}

%%%%%%%%%%%%%%%%%%%%%%%%%%%%%%%%%%%%%%%%%%%%%%%%%%%%%%%%%%%%
\section*{Preliminaries (Definitions and Lemmas)}
\begin{lemma}[Descent Lemma]\label{lem:descent}
If $f$ is $L$‑smooth, then for any $x,y\in C$,
\[
f(y)\le f(x)+\inner{\grad f(x)}{y-x}+ \frac{L}{2}\|y-x\|^{2}. 
\]
\end{lemma}
\begin{proof}
Standard consequence of the integral form of the gradient and the Lipschitz bound; see Assumption \ref{ass:smooth}. \hfill $\square$
\end{proof}
\begin{lemma}[Gap Upper Bound]\label{lem:gap}
For any $x\in C$ and $s\in C$,
\[
\gap(x)=\inner{\grad f(x)}{x-s}\ge 0,
\]
and for any $y\in C$,
\[
\inner{\grad f(x)}{y-x}\le \gap(x). 
\]
\end{lemma}
\begin{proof}
By definition $s=\arg\min_{u\in C}\inner{\grad f(x)}{u}$, therefore for all $y\in C$,
$\inner{\grad f(x)}{s}\le\inner{\grad f(x)}{y}$.
Rearranging gives the claim. \hfill $\square$
\end{proof}

%%%%%%%%%%%%%%%%%%%%%%%%%%%%%%%%%%%%%%%%%%%%%%%%%%%%%%%%%%%%
\section*{Main Proof (Step‑by‑Step Derivation)}
We prove Theorem \ref{thm:convergence}. Throughout the proof $k$ denotes the current iteration index.

\noindent\textbf{[Step 1] (Apply Descent Lemma to the update).}  
Using Lemma \ref{lem:descent} with $x=x_k$ and $y=x_{k+1}=(1-\gamma_k)x_k+\gamma_k s_k$,
\[
f(x_{k+1})\le f(x_k)+\inner{\grad f(x_k)}{x_{k+1}-x_k}
                +\frac{L}{2}\|x_{k+1}-x_k\|^{2}. \tag{5}
\]

\noindent\textbf{[Step 2] (Express the displacement).}  
Since $x_{k+1}-x_k=\gamma_k(s_k-x_k)$,
\[
\inner{\grad f(x_k)}{x_{k+1}-x_k}= \gamma_k\inner{\grad f(x_k)}{s_k-x_k}
= -\gamma_k\gap(x_k). \tag{6}
\]

\noindent\textbf{[Step 3] (Bound the squared norm).}  
Because $s_k,x_k\in C$, by definition of the diameter $D$,
\[
\|s_k-x_k\|\le D\;\Longrightarrow\;
\|x_{k+1}-x_k\|^{2}= \gamma_k^{2}\|s_k-x_k\|^{2}\le \gamma_k^{2} D^{2}. \tag{7}
\]

\noindent\textbf{[Step 4] (Combine (5)–(7)).}  
Substituting (6) and (7) into (5) yields
\[
f(x_{k+1})\le f(x_k)-\gamma_k\gap(x_k)+\frac{L}{2}\gamma_k^{2}D^{2}. \tag{8}
\]

\noindent\textbf{[Step 5] (Rearrange to isolate the gap).}  
Moving terms,
\[
\gamma_k\gap(x_k)\le f(x_k)-f(x_{k+1})+\frac{L}{2}\gamma_k^{2}D^{2}. \tag{9}
\]

\noindent\textbf{[Step 6] (Sum over $k=0,\dots,K-1$).}  
Summing (9) for $k=0$ to $K-1$,
\[
\sum_{k=0}^{K-1}\gamma_k\gap(x_k)
\le f(x_0)-f(x_K)+\frac{L D^{2}}{2}\sum_{k=0}^{K-1}\gamma_k^{2}. \tag{10}
\]

\noindent\textbf{[Step 7] (Lower bound the left‑hand side).}  
Because $\gamma_k\le 1$, we have $\gamma_k\gap(x_k)\ge \gamma_{\min}\gap(x_k)$ where
$\gamma_{\min}:=\min_{0\le k<K}\gamma_k$. For the chosen schedule $\gamma_k=2/(k+2)$ we have
$\gamma_k\ge \frac{2}{K+1}$ for all $k< K$, thus
\[
\sum_{k=0}^{K-1}\gamma_k\gap(x_k) \ge \frac{2}{K+1}\sum_{k=0}^{K-1}\gap(x_k). \tag{11}
\]

\noindent\textbf{[Step 8] (Upper bound the sum of squares).}  
For $\gamma_k=2/(k+2)$,
\[
\gamma_k^{2}= \frac{4}{(k+2)^{2}}\le \frac{4}{(k+1)(k+2)}.
\]
Hence
\[
\sum_{k=0}^{K-1}\gamma_k^{2}
\le 4\sum_{k=0}^{K-1}\Big(\frac{1}{k+1}-\frac{1}{k+2}\Big)
=4\Big(1-\frac{1}{K+1}\Big)\le 4 . \tag{12}
\]

\noindent\textbf{[Step 9] (Insert (11) and (12) into (10)).}  
Using $f(x_K)\ge f^\star$,
\[
\frac{2}{K+1}\sum_{k=0}^{K-1}\gap(x_k)
\le f(x_0)-f^\star+\frac{L D^{2}}{2}\cdot 4
= f(x_0)-f^\star+2L D^{2}. \tag{13}
\]

\noindent\textbf{[Step 10] (Derive the bound for the best gap).}  
Let $k^\star:=\arg\min_{0\le k<K}\gap(x_k)$. Then
\[
(K)\,\gap(x_{k^\star}) \le \sum_{k=0}^{K-1}\gap(x_k). \tag{14}
\]
Combining (13) and (14) gives
\[
\gap(x_{k^\star})
\le \frac{K+1}{2K}\big(f(x_0)-f^\star+2L D^{2}\big)
\le \frac{2\big(f(x_0)-f^\star\big)}{K}+ \frac{L D^{2}}{K}, \tag{15}
\]
which is exactly inequality \eqref{eq:4}. This completes the proof of the first part of Theorem \ref{thm:convergence}. \hfill $\square$

\noindent\textbf{[Step 11] (Average‑gap rate).}  
Dividing (13) by $K$ yields
\[
\frac{1}{K}\sum_{k=0}^{K-1}\gap(x_k)
\le \frac{f(x_0)-f^\star}{K}+ \frac{2L D^{2}}{K}
= \mathcal O\!\Big(\frac{1}{K}\Big).
\]
Applying Jensen’s inequality to the non‑negative sequence $\{\gap(x_k)\}$ gives the
$\mathcal O(1/\sqrt{K})$ bound stated in the theorem (standard conversion from
$1/K$ bound on the average to $1/\sqrt{K}$ bound on the minimum). \hfill $\square$

%%%%%%%%%%%%%%%%%%%%%%%%%%%%%%%%%%%%%%%%%%%%%%%%%%%%%%%%%%%%
\section*{Rate Derivation (With Constants)}
Collecting the constants appearing in \eqref{eq:4} we obtain:
\[
\boxed{\;
\min_{0\le k<K}\gap(x_k)
\;\le\;
\underbrace{\frac{2\big(f(x_0)-f^\star\big)}{K}}_{\text{initial suboptimality term}}
\;+\;
\underbrace{\frac{L D^{2}}{K}}_{\text{smoothness \& geometry term}}
\;}
\]
Hence, to achieve a target gap $\varepsilon>0$, it suffices to run at most
\[
K\ge \frac{2\big(f(x_0)-f^\star\big)+L D^{2}}{\varepsilon}
\]
iterations.

%%%%%%%%%%%%%%%%%%%%%%%%%%%%%%%%%%%%%%%%%%%%%%%%%%%%%%%%%%%%
\section*{Special Cases}
\begin{enumerate}[leftmargin=*]
\item \textbf{Convex $f$}. When $f$ is convex, the Frank–Wolfe gap upper‑bounds the primal gap:
\[
f(x_k)-f^\star \le \gap(x_k),
\]
so inequality \eqref{eq:4} directly yields the classic $\mathcal O(1/k)$ convergence of the objective value.

\item \textbf{Exact Line‑Search}. If $\gamma_k$ is chosen by solving
\[
\gamma_k = \arg\min_{\gamma\in[0,1]} f\big((1-\gamma)x_k+\gamma s_k\big),
\]
the term $\frac{L}{2}\gamma_k^{2}D^{2}$ in (8) is replaced by a non‑positive quantity,
leading to a tighter bound
\[
\gamma_k\gap(x_k)\le f(x_k)-f(x_{k+1}),
\]
and consequently $\gap(x_{k^\star})\le \frac{f(x_0)-f^\star}{K}$.

\item \textbf{Polytope $C$}. By Carathéodory’s theorem each iterate $x_k$ is a convex combination of at most $d+1$ extreme points of $C$, offering a sparsity guarantee absent in projected methods.

\end{enumerate}

%%%%%%%%%%%%%%%%%%%%%%%%%%%%%%%%%%%%%%%%%%%%%%%%%%%%%%%%%%%%
\section*{Discussion}
The proof above demonstrates that the Frank–Wolfe algorithm, despite its simplicity
and reliance only on a linear minimization oracle, attains the same worst‑case
first‑order convergence rate as projected gradient descent for smooth non‑convex
functions. The essential ingredients are the smoothness of $f$ (Descent Lemma) and
the compactness of the feasible region (finite diameter). The bound is
\emph{dimension‑free}—the constants involve only $L$, $D$, and the initial
suboptimality, not the ambient dimension $d$. This makes Frank–Wolfe particularly
attractive for large‑scale structured domains where projections are costly but
linear minimization is cheap.

\end{document}